\documentclass[10pt,a4paper]{amsart}
\usepackage[margin=19mm]{geometry}
\usepackage{amsmath,amssymb,mathtools}
\usepackage[scr=rsfs,cal=euler]{mathalfa}
\usepackage{xfrac}
\usepackage[unicode,hidelinks,bookmarks=false]{hyperref}
\newcommand{\ie}{i.e.\ }
\newcommand{\cf}{cf.\ }
\newcommand{\resp}{respectively\ }
\newcommand{\Subsection}{\subsection}
\providecommand{\nobreakdash}{\nobreak\hbox{-}}
\setlength{\emergencystretch}{2em}
\multlinegap0pt
\usepackage{mathtools}
\usepackage{amssymb}
\usepackage{xcolor}
\usepackage{bbm}
\usepackage{tikz}
\usepackage{dsfont}
\usepackage[shortlabels]{enumitem}
\usepackage{caption}
\newtheorem{theorem}{Theorem}
\newtheorem{claim}{Claim}
\newtheorem{proofclaim}{Proofclaim}
\usepackage{cleveref}
\crefname{theorem}{Theorem}{Theorems}
\crefname{claim}{Claim}{Claims}
\crefname{proofclaim}{Proofclaim}{Proofclaims}
\newcommand{\abs}[1]{\left\lvert#1\right\rvert}
\newcommand{\al}{\alpha}
\newcommand{\be}{\beta}
\newcommand{\ce}{\gamma}
\newcommand{\subs}{\subseteq}
\renewcommand{\subset}{\subseteq}
\title{Spanning Components and Surfaces Under Minimum Vertex Degree}
\author{Jack Allsop, Ander Lamaison, Richard Lang, Silas Rathke}
\date{Source: arXiv:2512.24242v1 (CC BY 4.0)}
\begin{document}
\maketitle

\noindent\emph{Source and licence.} Spanning Components and Surfaces Under Minimum Vertex Degree. Version arXiv:2512.24242v1 (CC BY 4.0), distributed by arXiv under the Creative Commons Attribution 4.0 International licence, http://creativecommons.org/licenses/by/4.0/, as recorded in the arXiv licence metadata for this exact version and shown on the abstract page https://arxiv.org/abs/2512.24242v1. Full licence notice: \texttt{licenses/arxiv-2512.24242-CC-BY-4.0.txt}.\par\smallskip

\noindent\textit{Editorial scope.} Counted unit: the article's theorem thm:spanning-component (source0.tex lines 610--612). The article's complete original proof of it is delivered with it (source0.tex lines 709--868, one \texttt{proof} environment of 160 lines, which the article places away from the statement). No mathematics is written, repaired or re-derived: the statement and the proof are the article's own lines, in the article's own notation, and no retained line is substituted or re-flowed.  No further source-local context is needed: every cross-reference of every retained line resolves inside the two delivered spans.  Nothing was omitted from any retained span, so the package carries no omission record.  No retained line contains a citation, so the wrapper carries no bibliography and no bibliography entry.  No retained line contains a figure environment, an image-inclusion command or drawing code, so no image is delivered and the package declares no asset dependency.  Editorial formatting only, in addition: an AMS \texttt{amsart} preamble that restates the article's own declarations and macros used by the retained lines, namely the environments \texttt{theorem}, \texttt{claim}, \texttt{proofclaim} on separate counters, together with the standard packages those lines need.  The retained environments print in this document as Theorem 1, Claim 1, Proofclaim 1 in order of appearance, whereas the article prints the same items with the numbering of its own full document.  The complete native source of the article as distributed in the arXiv e-print for this exact version is bundled unchanged as \texttt{sources/191868/source0.tex}.
	\begin{theorem}\label{thm:spanning-component}
		Every $3$-graph $G$ on $n$ vertices with $\delta_1(G) \geq \tfrac{1}{2}\binom{n-1}{2}+1$ contains a spanning component.
	\end{theorem}

	\begin{proof}[Proof of \cref{thm:spanning-component}]
		Let $G$ be a $3$-graph on $n$ vertices with $\delta_1(G) \geq \tfrac{1}{2}\binom{n-1}{2}+1$.
		Consider a vertex $x \in V(G)$, and note that $L(x)$ contains at most $n-1$ non-isolated vertices and at least $\tfrac{1}{2} \binom{n-1}{2} + 1$ edges since $\delta_1(G) \geq \tfrac{1}{2} \binom{n-1}{2} + 1$. Therefore, $L(x)$ contains a vertex of degree at least $(n-1)/2$. 
		So in particular, $L(x)$ has a unique component of largest order, which we denote by $C_x$.
		Let $\hat C_x \subset G$ contain the edges $xyz$ with $yz \in C_x$.
		Since $C_x$ is connected in $L(x)$, it follows that $\hat C_x$ is connected in $G$.
		So $\hat C_x$ is contained in a (tight) component of $G$, which we denote by $C_x^\ast$.
		
		For the sake of contradiction, let us assume that $G$ does not contain a spanning component.
		Let $c$ be the number of distinct components $C_v^\ast$ over all $v \in V(G)$.
		Evidently, $c \geq 2$.
		In the following, we distinguish between the cases $c=2$, $c=3$, and $c\geq 4$.
		
		\medskip
		\textbf{Case 1: $c=2$.}
		Fix a vertex $x$ that does not appear in all $V(C_v^\ast)$. By assumption, $C_x^\ast$ is not a spanning component, so there is a vertex $y\not\in V(C_x^\ast)$. Since $c=2$, every $C_v^\ast$ is either $C_x^\ast$ or $C_y^\ast$. In particular, $x\not\in V(C_y^\ast)$.
		
		Let $V'\coloneqq V(G)\backslash\{x,y\}$, $L'(x)\coloneq L(x)[V']$, and $L'(y)\coloneqq L(y)[V']$. 
		Since $\delta_1(G) \geq \tfrac{1}{2} \binom{n-1}{2} + 1$, the link graph $L(x)$ has at least $\frac{1}{2}\binom{n-1}{2}+1$ edges. Fewer than $\frac{n}{2}$ of them can be incident to $y$ since, otherwise, $y$ would be in the largest component $C_x$ of $L(x)$. Thus, $e(L'(x))\ge \frac{1}{2}\binom{n-1}{2}+1-\frac{n}{2}$. The same is true for $L'(y)$. This implies \begin{align*}
			e\big(L'(x)\big)+e\big(L'(y)\big)&>2\cdot\!\left(\frac{1}{2}\binom{n-1}{2}+1-\frac{n}{2}\right)
			%			=\frac{n^2-5n+6}{2}
			=\binom{n-2}{2}\\&=e\big(\overline{L'(x)}\cap\overline{L'(y)}\big)+e\big(L'(x)\big)+e\big(L'(y)\big)-e\big(L'(x)\cap L'(y)\big)\,.
		\end{align*}
		We can conclude 
		\begin{equation}\label{eq: more in cap than the complement}
			e\big(L'(x)\cap L'(y)\big)>e\big(\overline{L'(x)}\cap\overline{L'(y)}\big)\,.
		\end{equation}
		
		We define
		\begin{align*}A &\coloneqq \big\{v\in V'\colon v\in V(C_x), v\in V(C_y)\big\}\\
			B &\coloneqq \big\{v\in V'\colon v\in V(C_x), v\not\in V(C_y)\big\}\\
			D &\coloneqq \big\{v\in V'\colon v\not\in V(C_x), v\not\in V(C_y)\big\}\,.
		\end{align*}
		We denote the sizes of the sets $A$, $B$, and $D$ by $a$, $b$, and $d$, respectively.
		Because of \[b+d=n-2-\abs{V(C_y)}\le \frac{n}{2}-2\le \abs{V(C_x)}-2<\abs{V(C_x)}=a+b,\]
		we have $d< a$.
		
		By the assumption on the minimum degree, there is an edge $\{u,v\}\in E\big(L'(x)\cap L'(y)\big)$, that is,  $\{u,v,x\},\{u,v,y\}\in E(G)$.
		We show now that $u,v\in D$. Suppose $u\in V(C_x)$. Then there is a $w$ such that $\{u,w,x\}\in C_x^\ast$, but then we also have that $\{u,v,x\}$ and $\{u,v,y\}$ are in $C_x^\ast$ which contradicts the assumption that $y\not\in V(C_x^\ast)$. Thus, $u\not\in V(C_x)$ and in the same way one can show that $u\not\in V(C_y)$. Hence, $u\in D$ and, similarly, $v\in D$. In particular, this implies that for any $e\in E\big(L'(x)\cap L'(y)\big)$, its endpoints must both lie in $D$. Thus, $e\big(L'(x)\cap L'(y)\big)\le \binom{d}{2}\le d^2$. This also implies $d\ge 1$.
		
		Next, consider a pair $\{u,v\}$ with $u\in A$ and $v\in D$. We cannot have $\{u,v\}\in L'(x)$, since then $u\in A\subs V(C_x)$ would imply that $v\in V(C_x)$ contradicting $v\in D$. Similarly, we cannot have $\{u,v\}\in L'(y)$. Thus, $\{u,v\}\in E\big(\overline{L'(x)}\cap\overline{L'(y)}\big)$ and
		$e\big(\overline{L'(x)}\cap\overline{L'(y)}\big)\ge ad$. Therefore, \eqref{eq: more in cap than the complement} implies
		\[ad\le e\big(\overline{L'(x)}\cap\overline{L'(y)}\big)<e\big(L'(x)\cap L'(y)\big)\le d^2\,,\]
		that is, $a<d$ which contradicts our previous result $d< a$.
		
		\medskip
		\textbf{Case 2: $c=3$.}
		Fix vertices $x$, $y$, and $z$ such that $C_x^\ast$, $C_y^\ast$, and $C_z^\ast$ are pairwise distinct. Let $C_x'$ be the graph on vertex set $V(G)$ and edge set $E(C_x)\cup\{ux\colon u\in V(C_x)\}$. We define $C_y'$ and $C_z'$ in a similar way. Note that every pair in $\binom{V(G)}{2}$ can be in at most one of $E(C'_x)$, $E(C'_y)$, and $E(C'_z)$ as, otherwise, $C_x^\ast$, $C_y^\ast$, and $C_z^\ast$ would not be distinct. In particular $e(C'_x)+e(C'_y)+e(C_z')\le \binom{n}{2}$.
		
		Denote $A \coloneq V(C_x)$, $B\coloneq V(C_y)$, and $C \coloneq V(C_z)$.
		Set $\al \coloneq\frac{\abs A}{n-1}$, $\be\coloneq\frac{\abs B}{n-1}$, and $\ce\coloneq\frac{\abs C}{n-1}$. Note that $\al,\be,\ce\in[\frac12,1]$. In the following, we will often use that for $\delta\in[0,1]$, we have
		\begin{align}\label{eq:pull out}
			\binom{\delta(n-1)}{2}=\delta^2\binom{n-1}{2}-\frac{(1-\delta)\delta}{2}(n-1) \,.
		\end{align}
		
		Now observe that at most $\binom{n-1-\al(n-1)}{2} =\binom{(1-\al)(n-1)}{2}= (1-\alpha)^2\binom{n-1}{2}-\frac{\alpha(1-\alpha)}2(n-1)$ edges of $L(x)$ are not in $C_x$.
		This, together with the fact that $e(L(x))\ge \frac{1}{2}\binom{n-1}{2}+1$, implies $e(C_x) > \big(1/2 - (1-\al)^2\big)\binom{n-1}{2}+\frac{\alpha(1-\alpha)}2(n-1)$ and, ignoring the second summand,  $e(C'_x)\ge \frac{n-1}{2}+e(C_x)>\frac{n-1}{2}+ \big(1/2 - (1-\al)^2\big)\binom{n-1}{2}$. Similarly, $e(C'_y)>\frac{n-1}{2}+ \big(1/2 - (1-\be)^2\big)\binom{n-1}{2}$ and $e(C'_z)>\frac{n-1}{2}+ \big(1/2 - (1-\ce)^2\big)\binom{n-1}{2}$.
		Therefore,
		\begin{align}
			\binom n2 &\ge e(C_x')+e(C_y')+e(C_z')\notag\\&>\frac{3}{2}(n-1)+\left(\frac{3}{2}-\big(1-\al\big)^2-\big(1-\be\big)^2-\big(1-\ce\big)^2\right)\binom{n-1}{2}\notag\\\label{equ:Cx-cup-Cy-cup-Cz}
			&>n-1+\left(\big(2\al-\al^2\big)+\big(2\be-\be^2\big)+\big(2\ce-\ce^2\big)-\frac{3}{2}\right)\binom{n-1}{2}\,.
		\end{align}
		Suppose first that $\al+\be+\ce\geq 2$.
		Then the right side is minimised, for example, when $\al=1/2$, $\be=1/2$, and $\ce=1$ since $f(a)=2a-a^2$ is both concave and monotone increasing on $\big[\frac{1}{2},1\big]$. However, in this case the right side of \eqref{equ:Cx-cup-Cy-cup-Cz} equals 
		\[n-1+\binom{n-1}{2}=\binom{n}{2}\,,\]
		a contradiction. So we can assume that $\al+\be+\ce\leq 2$.
		
		Recall that $C_x$, $C_y$, and $C_z$ are edge-disjoint.
		Together with the arguments after \eqref{eq:pull out}, we obtain
		\begin{align}
			\label{equ:Cx-cup-Cy-cup-Cz-noprime}
			e(C_x \cup C_y \cup C_z) &= e(C_x)+e(C_y)+e(C_z) \notag \\&> \left(\big(2\al-\al^2\big)+\big(2\be-\be^2\big)+\big(2\ce-\ce^2\big)-\frac{3}{2}\right)\binom{n-1}{2}\notag\\
			&\quad+\left(\frac{\alpha(1-\alpha)}{2}+\frac{\beta(1-\beta)}{2}+\frac{\gamma(1-\gamma)}{2}\right)(n-1)\notag\\
			&=f_1(\alpha,\beta,\gamma)\binom{n-1}{2}+f_2(\alpha,\beta,\gamma)(n-1)\,,
		\end{align}
		where $f_1(\alpha,\beta,\gamma)\coloneq \big(2\al-\al^2\big)+\big(2\be-\be^2\big)+\big(2\ce-\ce^2\big)-\frac{3}{2}$ and $f_2(\alpha,\beta,\gamma)\coloneq \frac{\alpha(1-\alpha)}{2}+\frac{\beta(1-\beta)}{2}+\frac{\gamma(1-\gamma)}{2}$.
		
		Let us now find an upper bound for the left side of \eqref{equ:Cx-cup-Cy-cup-Cz-noprime}.
		By the principle of inclusion and exclusion, we obtain that
		\begin{align}
			e(C_x \cup C_y \cup C_z) &\leq \binom{|A|}{2} + \binom{|B|}{2} + \binom{|C|}{2} +  \binom{|A \cap B \cap C|}{2}  \nonumber \\ 
			&\quad - \binom{|A \cap B|}{2} - \binom{|B \cap C|}{2} - \binom{|C \cap A|}{2} \,. \label{equ:AcupBcupC}
		\end{align}
		We denote the right side of \eqref{equ:AcupBcupC} by $g(A,B,C)$. We claim that there are sets $A',B',C'$ with $\abs{A'}=\abs{A}$, $\abs{B'}=\abs{B}$, $\abs{C'}=\abs{C}$, where $g(A,B,C)\le g(A',B',C')$ and $A'\cap B'\cap C'=\emptyset$. Indeed, if already $A\cap B\cap C=\emptyset$, then there is nothing to show. Suppose that there is a $v \in A \cap B \cap C$.
		Since $\abs{A} + \abs{B} + \abs{C} =(\al+\be+\ce)(n-1)< 2n$, there is an element $u$ that is in at most one of the sets $A$, $B$, or $C$, say $u \in A$.
		We may then delete~$v$ from $B$ and add $u$ instead.
		It follows that the terms $|B \cap C|$ and $|A \cap B \cap C|$ decrease each by one, while $|A|$, $|B|$ and $|C|$ stay the same.
		Since $|B \cap C| \geq |A \cap B \cap C|$, this implies that $g(A,B,C)$ does not decrease when we make this change. Since $\abs{A\cap B\cap C}$ decreases for each such change, eventually, we end up with the desired sets $A',B',C'$.
		
		Note that $\abs{A' \cap B'} \geq \abs{A'} + \abs{B'} -(n-1)$ and similar inequalities hold for $\abs{A' \cap C'}$ and $\abs{B' \cap C'}$. Using this, we may bound the right side of \eqref{equ:AcupBcupC} from above by
		\begin{align*}
			g(A',B',C') &\le \binom{\al(n-1)}{2}+\binom{\be(n-1)}{2}+\binom{\ce(n-1)}{2}
			\\&\quad-\binom{(\al+\be-1)(n-1)}{2} 
			\\&\quad-\binom{(\be+\ce-1)(n-1)}{2} 
			\\&\quad-\binom{(\ce+\al-1)(n-1)}{2}\\
			&\overset{(\ref{eq:pull out})}= \left(\al^2+\be^2+\ce^2-(\al+\be-1)^2-(\be+\ce-1)^2+(\ce+\al-1)^2\right)\binom{n-1}{2} \\
			& \quad+\left(-\frac{(1-\alpha)\alpha}{2}-\frac{(1-\beta)\beta}{2}-\frac{(1-\gamma)\gamma}{2}\right)(n-1)\\
			&\quad +  \frac{(2-\alpha-\beta)(\alpha+\beta-1)}{2} (n-1)
			\\&\quad+\frac{(2-\beta-\gamma)(\beta+\gamma-1)}{2} (n-1)
			\\&\quad+\frac{(2-\gamma-\alpha)(\gamma+\alpha-1)}{2} (n-1)\\
			&=f_3(\alpha,\beta,\gamma)\binom{n-1}{2}+f_4(\alpha,\beta,\gamma)(n-1)\,,
		\end{align*}
		where 
		$$f_3(\alpha,\beta,\gamma)\coloneq\al^2+\be^2+\ce^2-(\al+\be-1)^2-(\be+\ce-1)^2+(\ce+\al-1)^2$$ 
		and 
		\begin{align*}
			f_4(\alpha,\beta,\gamma) & \coloneq -\frac{(1-\alpha)\alpha}{2}-\frac{(1-\beta)\beta}{2}-\frac{(1-\gamma)\gamma}{2}
			\\&\quad+\frac{(2-\alpha-\beta)(\alpha+\beta-1)}{2}
			\\&\quad+\frac{(2-\beta-\gamma)(\beta+\gamma-1)}{2}
			\\&\quad+\frac{(2-\gamma-\alpha)(\gamma+\alpha-1)}{2}\,.
		\end{align*}
		Combining this with the right side of \eqref{equ:Cx-cup-Cy-cup-Cz-noprime} gives
		\begin{align}\label{equ:final}
			f_1(\alpha,\beta,\gamma)\binom{n-1}{2}+f_2(\alpha,\beta,\gamma)(n-1)<f_3(\alpha,\beta,\gamma)\binom{n-1}{2}+f_4(\alpha,\beta,\gamma)(n-1)\,.
		\end{align}
		We will derive a contradiction by showing the following two claims.
		
		\begin{claim}
			We have $f_1(\alpha,\beta,\gamma)\ge f_3(\alpha,\beta,\gamma)$.
		\end{claim}
		\begin{proofclaim}
			For sake of contradiction, suppose that $f_1(\alpha,\beta,\gamma) < f_3(\alpha,\beta,\gamma)$.
			Rearranging gives
			\begin{align*}
				0 & > 2(\al\be +\be\ce + \ce\al) - 2(\al+\be+\ce) + 3/2 \\ & = 2(\al-1/2)(\be-1/2) + 2(\be-1/2)(\ce-1/2) + 2(\al-1/2)(\ce-1/2)   \,.
			\end{align*}
			However, this contradicts the fact that $\alpha$, $\beta$, and $\gamma$ are each at least $1/2$.
		\end{proofclaim}
		
		\begin{claim}
			We have  $f_2(\alpha,\beta,\gamma)\ge f_4(\alpha,\beta,\gamma)$.
		\end{claim} 
		\begin{proofclaim}
			To begin, observe that 
			\begin{align*}
				0&\le 3-\frac{3}{2}(\al+\be+\ce) \,,
			\end{align*}
			because $\al+\be+\ce\le2$.
			Moreover,
			\begin{align*}
				0&\le \al\be+\al\ce+\be\ce-\frac{1}{2}(\al+\be+\ce)\,,
			\end{align*}
			since $\al,\be,\ce\ge \frac{1}{2}$ implies $\al\be\ge \frac{1}{2}\al$, $\be\ce\ge\frac{1}{2}\be$, and $\ce\al\ge\frac{1}{2}\ce$.
			
			Now assume contrariwise that  $f_2(\alpha,\beta,\gamma) < f_4(\alpha,\beta,\gamma)$.
			%https://www.wolframalpha.com/input?i=a%281-a%29%2F2%2Bb%281-b%29%2F2%2Bc%281-c%29%2F2+%3E%3D+-%281-a%29a%2F2-%281-b%29b%2F2-%281-c%29c%2F2%2B%282-a-b%29%28a%2Bb-1%29%2F2%2B%282-b-c%29%28b%2Bc-1%29%2F2%2B%282-c-a%29%28c%2Ba-1%29%2F2
			Simplifying and rearranging gives
			\[0 > 3+\al\be+\al\ce+\be\ce-2(\al+\be+\ce) \,.\]
			But this contradicts at least one of the above inequalities.
		\end{proofclaim}
		
		By combining these two claims, we obtain a contradiction to \eqref{equ:final}, which concludes the case when $c=3$.
		
		
		
		\medskip
		\textbf{Case 3: $c\geq4$.} Following the same steps as in case 2, we can derive the analogue of~\eqref{equ:Cx-cup-Cy-cup-Cz}, with four components instead of three. Thus, there are $\alpha,\beta,\gamma,\eta\in \big[\frac{1}{2},1\big]$ such that
		\[\binom{n}{2}>n-1+\left(\big(2\al-\al^2\big)+\big(2\be-\be^2\big)+\big(2\ce-\ce^2\big)+\big(2\eta-\eta^2\big)-\frac{4}{2}\right)\binom{n-1}{2} \,.\]
		But the right side is minimised for $\alpha=\beta=\gamma=\eta=\frac{1}{2}$, where it is $n-1+\binom{n-1}{2}=\binom{n}{2}$, which is absurd!
	\end{proof}

\end{document}
